\begin{textsample}{POS Dim 4 – ma – Score 13.00 – ma/ma\_ef\_59.txt}  \label{ex:f4_pos_055}
3.5.1 ENSINO \textbf{RELIGIOSO} O Ensino \textbf{Religioso}, como componente curricular da Educação Básica, expressa clareza quanto à sua \textbf{identidade} e natureza subjetiva e de reflexão transcendental, contudo é composto de múltiplos objetos de estudo, abordados em conteúdos diversos, com uma metodologia variável de acordo com o conteúdo a ser trabalhado, bem como uma avaliação contínua desenvolvida numa perspectiva dinâmica salientando a ludicidade, progredindo em complexidade, a partir das fases de análise, síntese e crítica, ao longo do processo de formação do estudante.
Do ponto de vista da prática, do fazer pedagógico e docente, o Ensino \textbf{Religioso} deve estimular no estudante “ a pesquisa e o diálogo como princípios mediadores e articuladores dos processos de observação, identificação, análise, apropriação e ressignificação dos saberes visando o desenvolvimento das competências específicas ” ( BRASIL, 2017:434 ).
Entre as 10 competências gerais estabelecidas pela BNCC, algumas são extremamente pertinentes à perspectiva de um Ensino \textbf{Religioso} que tem como base pressupostos éticos, filosóficos e científicos.
São elas : a que permite valorizar os conhecimentos construídos ao longo do processo histórico de construção da humanidade no sentido de tentar explicar a realidade, quer seja física, espiritual, social, cultural e digital ; a que estabelece a fruição das diversas manifestações artísticas e culturais, locais, regionais, nacionais e mundiais, bem como as diversidades de saberes e vivências culturais ; a que incentiva o exercício da cidadania com consciência crítica, liberdade, autonomia e responsabilidade e as que propõem a elaboração de argumentos com base na defesa dos direitos humanos, com posições éticas, no sentido do cuidado de si, dos outros e do planeta, exercitando a empatia, o diálogo, a resolução de conflitos e a cooperação com valorização da diversidade de indivíduos e de grupos sociais e suas \textbf{identidades} e um agir pessoal e coletivo com base em princípios éticos, democráticos, inclusivos, sustentáveis e solitários, colaborando desta forma no desenvolvimento de uma sociedade mais justa, equânime e democrática ( BRASIL, 2017 ).
Pensar sobre as competências específicas para o Ensino \textbf{Religioso} é refletir sobre o currículo.
O currículo é intencional, ou seja, tem intencionalidade.
Uma dessas intencionalidades é a forma-ção da personalidade do indivíduo e deve considerar, evitando os conflitos, as reformas educativas e curriculares em seus quatro níveis, a saber : federal ( LDB, DCNs e BNCC ), estadual ( referenciais curriculares ), escola ( PPP ) e a sala de aula ( metodologias das áreas de conhecimento ).
Neste sentido, convém ressaltar a concepção de Brandão ( 2005:78 ) : > O artigo 26 da LDB trata dos conteúdos curriculares do Ensino Fundamental e Médio, os quais devem ter uma base nacional comum, a ser complementada, em cada sistema de ensino e estabelecimento escolar, por uma parte diversificada, exigida pelas características regionais e locais da sociedade, da cultura, da economia e a clientela.
Desta forma, a historicidade do estudante deve ser considerada no processo de apropriação do conhecimento, e o currículo deve assegurar uma proposta diversificada, respeitando as diferenças individuais e culturais de cada um e observando que a religião é uma prática coletiva.
Aprendemos com a Antropologia que a ideia de diversidade encontra -se intimamente associada à noção do outro.
O Ensino \textbf{Religioso} deve dialogar com os valores morais e éticos numa convivência baseada no respeito e na busca pela paz, considerando a religiosidade de cada um dentro e fora da família, considerando a multiplicidade de crenças e tradições \textbf{religiosas}.
Conforme as Diretrizes Curriculares da rede estadual de ensino do ( MARANHÃO, 2014:75 ) : > O ser humano se desenvolve na medida em que se expressa e se relaciona.
Da mesma forma, a religiosidade torna -se efetiva e se desenvolve pela expressividade, comunicabilidade e linguagem.
O dinamismo da religiosidade ganha forma, ritmo e intensidade no fenômeno \textbf{religioso}.
Desta maneira, tratar metodologicamente o fenômeno religioso/experiência \textbf{religiosa} como fonte de tensão entre aquilo que é instituído e \textbf{hierarquizado}, frente à experiência propriamente dita do sujeito e do outro, pode ser um caminho de emergência verdadeira de conhecimento, de um conhecimento propriamente \textbf{religioso}.
Por esta razão, o Ensino \textbf{Religioso} como componente curricular apresenta -se com um forte elemento filosófico, no sentido de estabelecer um movimento contínuo de construção e desconstrução das verdades, valores e significados produzidos historicamente pelas práticas humanas.
Silva ( 2004:2-6 ) destaca que : > Em nenhum período da história houve uma única religião em todo o mundo, como também nunca foram dominantes as atitudes de tolerância no passado da história das religiões.
A associação entre Estado e Igreja é uma dessas formas de \textbf{intolerância}, não deixando, por isso mesmo, uma boa lembrança.
A imposição de uma fé como oficial e a consequente \textbf{exclusão} das outras ( inclusive com perseguições declaradas ) deixou seu rastro perverso no passado.
No presente, muitos conflitos continuam sendo alimentados a partir de convicções ou sob a justificativa de crença, como vemos no Oriente Médio ou na Irlanda. (... ) Nenhuma tradição \textbf{religiosa} é “ total ”, nem existe um status de favoritismo de religiões.
Conhecer o lugar onde estamos e onde os outros estão em relação à fé e às crenças leva -nos a desenvolver um sentido de proporção no amplo campo das religiões, religiosidades, experiências \textbf{religiosas} - onde todos devem ser ouvidos e respeitados.
A diversidade se faz riqueza e deve conduzir à compreensão, respeito, admiração e atitudes pacificadoras.
Numa sociedade democrática e no \textbf{Estado} \textbf{laico}, na perspectiva de uma educação pública de formação do cidadão crítico e reflexivo, deve -se lutar sempre por um Ensino \textbf{Religioso} que não expresse os preceitos e dogmas de uma única ou qualquer religião, sendo entendido como catequese ou pregação apologética.
O Ensino \textbf{Religioso} que se propõe aqui tem como princípio fundamental a diversidade, pois somos diversos histórica, étnica e linguisticamente e da mesma forma somos diversos religiosamente.
O Ensino \textbf{Religioso}, na perspectiva aqui apresentada, não se ocupa com ensinamentos de cunho denominacional ou congregacional, focando sua função social na promoção da liberdade \textbf{religiosa} e nos direitos de aquisição de conhecimentos diversos, não contemplando controvérsias doutrinárias de uma única religião.
Ao ser qualificado como “ \textbf{religioso} ” submete -se a uma área específica de atuação que tem como destinatário o sujeito, \textbf{religioso} ou não, que indaga sobre as razões de ser \textbf{religioso} dentro ou fora da religião, a partir de dentro ou de fora do grupo \textbf{religioso}, ou em não se ter religião alguma ( FIGUEIREDO, 1995:41-51 ).
Sendo assim, ele deve fundamentar seus conteúdos no estudo das culturas e tradições \textbf{religiosas}, das diversas teologias, fontes de textos filosóficos, dos ritos, simbologias e significados culturais e \textbf{religiosos}, e da ética.
O ensino \textbf{religioso} visa subsidiar as reflexões a respeito do imanente, do transcendente e do ser humano em sua diversidade de modelos de fé e de vida \textbf{religiosa}.
Segundo Junqueira ( 2007:14 ), não é “ função do Ensino \textbf{Religioso} escolar, promover conversões, mas oportunizar ambiente favorável para a experiência do Transcendente, em vista de uma educação integral, atingindo as diversas dimensões da pessoa ”.
Assim, é essencial a valorização das \textbf{identidades}, experiências e cosmovisões, seja do próprio interlocutor ou daqueles que o cercam, respeitando as raízes culturais, nas mais diversas perspectivas interculturais.
A valorização das religiões \textbf{indígenas}, africanas, afro-brasileiras, judaica, cristã, islâmica, espírita, entre outras, e dos conhecimentos não \textbf{religiosos} – ateísmo, materialismo, entre outros –, deverá ser assumida como princípio para pesquisa e diálogo respeitoso e acolhedor, promovendo reais processos de análise, apropriação e ressignificação de saberes.
O Maranhão possui características muito próprias no que se refere à religiosidade.
As mais comuns são de um catolicismo \textbf{popular} ( como : rezas, orações e santos ), culturas \textbf{indígenas} ( divindades e rituais de cura ) e culturas africanas ( ritos e tambores ).
Segundo Barros ( 2007:1 ) : > O Maranhão, localizado a dois graus ao \textbf{sul} do Equador, na \textbf{fronteira} sociogeográfica entre a Amazônia e o \textbf{Nordeste} do Brasil, possui uma diversidade de práticas culturais e \textbf{religiosas}, o que se relaciona ao conjunto múltiplo de povos que formaram essa região e à heterogeneidade das interações entre eles estabelecidas.
Eram diversos os povos nativos que habitavam esse torrão quando da vinda dos primeiros europeus no século XVI.
A estrutura social da região foi ainda mais complexificada com a chegada massiva de africanos a partir do século XVIII, quando o Maranhão, assim como a Bahia, passou a constituir uma das áreas mais negras do Brasil, e continuou, do mesmo modo que a Amazônia, uma importante região \textbf{indígena}.
A cultura, o conhecimento e a religiosidade se expressam por meio da linguagem que se comunica com diversos segmentos da sociedade e, muitas vezes, definem valores e consequentemente estabelecem normas e comportamentos sociais, caracterizando assim uma sociedade.
Esses traços indeléveis de cada cultura local constituem o \textbf{território} maranhense, tornando -o único e diverso, pois as diferenças que compõem a sua unidade devem ser consideradas, sobremaneira, em um Ensino \textbf{Religioso} que, aqui, propõe -se a adotar “ a pesquisa e o diálogo como princípios mediadores e articuladores dos processos de observação e identificação, análise, apropriação e ressignificação de saberes, visando o desenvolvimento de competências especificas ” ( BNCC, 2017:434 ).
Sendo assim, o diálogo inter-religioso e intercultural, baseado no fundamento ético do respeito à alteridade, é fundamental para o reconhecimento da diversidade \textbf{religiosa}, implicando responsabilidades mútuas para que haja uma boa convivência, enquanto princípio norteador de escolhas, atitudes e políticas de vida coletiva.
Pensar o Ensino \textbf{Religioso} como componente curricular representa, na escola, o rumo traçado para que se efetive uma prática pedagógica fundada no respeito às diferenças e diversidades, solidificando os valores de respeito e tolerância à crença de cada um e de cada cultura.
Por isso, deve estar atrelado às necessidades sociais do momento, para tornar concreto o objetivo a que se propõe a escola.
Deve ser dinâmico e altamente adaptável às circunstâncias sociais e suas exigências.
O componente curricular Ensino \textbf{Religioso} tem por finalidade debater a perspectiva do sagrado no ser humano, manifestado em circunstâncias históricas, culturais e \textbf{religiosas} diversas ao longo do tempo e do espaço.
Tendo em vista a mudança de paradigmas desenhada a partir dos pensamentos de alguns teóricos, tais como Einstein, Stanislau, Capra, Kuhn, Morin, Varela, Latour, entre outros, surgem novos olhares sobre o mundo e o conhecimento, abrindo espaço para as concepções sistêmica, holística e complexa, desencadeando uma nova cosmovisão.
Essas visões que aparecem numa perspectiva mais ampla, também na sociedade da informação e do conhecimento, começam a ser pensadas numa perspectiva teológica, possibilitando assim novas reflexões sobre as religiosidades, inclusive com o aparecimento de novas religiões e crenças ou de novas denominações de matrizes \textbf{religiosas} já existentes.

% matched lemmas: estado, exclusão, fronteira, hierarquizar, identidade, indígena, intolerância, laico, nordeste, popular, religioso, sul, território
\end{textsample}
